% New input section; independent successor to the 120-step construction.
\section{A coefficient normalization allowing shorter masks}
\label{sec:short-masks}

The 120-step system still pays for two powers that can be removed by
changing the coefficient encoding. Its first coefficient code has a
nonzero digit in position $L$, and the packing shifts the second code by
$q^2$. We construct a quartic representation whose digit in position $L$
is zero. Both coefficient codes can then be bounded by $q$, their packing
uses a shift by $q$, and the block widths can be reduced to
$q,q^2,q^5$. Their product is $q^8$; only $q^2,q^4,q^8$ need be computed.

The normalization is part of the admissible index, not an assumption
about an arbitrary polynomial. In particular, simply deleting the leading
digit of the old code would change the coefficient being tested. The
construction below makes that digit zero while preserving the represented
set.

\subsection{Making the leading coefficient digit vanish}

Start with the quadratic-system representation underlying the universal
pair $(58,4)$, and shift its witnesses to nonnegative coordinates. Write
its residuals as $F_i(X,z_1,\ldots,z_{58})$, with constant coefficients
$c_i$. Introduce a nonnegative witness $h$ and put
\[
 F_i^h=F_i-c_i+c_ih,
 \qquad
 S=\sum_i(F_i^h)^2+(Xh-X)^2.
\]
The polynomial $S$ has degree at most four and zero constant coefficient.
For $X>0$, $S=0$ forces $h=1$ and hence is equivalent to the original
system. Introduce a further nonnegative witness $\sigma$ and choose a
power of two $z\ge4$. The polynomial to be encoded is
\begin{equation}
 R=24S+z(\sigma-1).\label{eq:short-normalization}
\end{equation}
For nonnegative integral witnesses, $R=0$ is equivalent to $S=0$ and
$\sigma=1$: values $\sigma\ge2$ make $R$ positive, while $\sigma=0$
would require $24S=z$, impossible because $3$ does not divide a power of
two. Thus this remains a universal quartic representation for positive
inputs. It now uses $\nu=60$ witness coordinates, including $h,\sigma$.
These are coordinates encoded in $g$, not additional unknowns of the
final twenty-equation system.
The forced coordinate $\sigma=1$ ensures that the actual witness code
$g$ is positive.

For multiindices $I=(i_0,\ldots,i_\nu)$ of total degree at most four,
put
\[
 c_I=\frac{4!}{i_0!\cdots i_\nu!(4-\sum_j i_j)!},
 \qquad R=\sum_I c_IP_I X^{i_0}z_1^{i_1}\cdots z_\nu^{i_\nu}.
\]
Every $c_I$ divides 24. Consequently the contribution of $24S$ to every
$P_I$ is integral; the new linear coefficient of $\sigma$ contributes
$z/4$, also integral. Choose $z$ so large that the fixed contributions
from $24S$ have absolute value less than $z$. The variable $\sigma$ is
fresh and hence these contributions do not depend on $z$. We obtain
\begin{equation}
 P_{\boldsymbol0}=-z,
 \qquad -z<P_I<z\quad(I\ne\boldsymbol0).
 \label{eq:short-coefficients}
\end{equation}
The convention here encodes $R$ directly as $\sum c_IP_I$ times the
monomials. There is no further factor of 24 in the extracted coefficient.

Set
\[
 Z=2z,\qquad M=5^{60},\qquad L=5^{64}=625M,
 \qquad w(I)=i_0+i_1 5+\cdots+i_\nu5^\nu.
\]
Then $w(I)\le4M<L$, and distinct $I$ have distinct weights by uniqueness
of base-five digits. Define
\begin{align*}
 \ell_0(B)&=\sum_{j=1}^{\nu}B^{5^j},\\
 e_0(B)&=\sum_I(z+P_I)B^{L-w(I)},\\
 V&=e_0(Z)+\ell_0(Z)Z^L.
\end{align*}
The leading digit of $e_0$ vanishes by
\eqref{eq:short-coefficients}, so $\deg e_0\le L-1$. All its other
digits lie in $[0,Z)$; the two parts of the packed polynomial
$e_0(B)+\ell_0(B)B^L$ have disjoint positions.

Choose a power of two $H$ with
\begin{equation}
 H>\max\{2Z^{2L+1},\ Z(\nu+2)^4,\ 3L,\ 64\}.
 \label{eq:short-H}
\end{equation}
An index $(Z,V,H)$ is now called admissible if constructed by these rules
from a quadratic-system representation. Its three components remain
fixed parameters. The complete parameter list is $x,Z,V,H$, as in the
preceding packed system, although this is a different class of admissible
indices.

\subsection{The shortened system and its certificate}

Keep the 32 positive unknowns of $\Sigma_{\rm pack}$ and the
abbreviations $B=Hb^4$, $C=1+xB+g$. Replace the coding equations by
\begin{align}
 b&=x+\beta,& e+\ell C^4+\alpha&=q,\label{eq:short-bound}\\
 \lambda+q^2&=1+\lambda B,&\theta+Z&=B,\label{eq:short-geometric}\\
 e+\ell q&=V+t\theta,&n&=q^8.\label{eq:short-packed}
\end{align}
Use the following abbreviations in the equation defining $r$:
\begin{align*}
 S_3&=(2e-Z\lambda)C^4+B\lambda(1+q^2),\\
 S&=g+q(e+\ell q+q^2S_3),\\
 T+1&=q-(b-1)\ell+\theta\lambda q+(B-2)q^4,\\
 r&=S(n^2-n)+(T+1)(n^2-1).
\end{align*}
Retain the remaining Pell equations, including $c=\kappa+\varphi$,
and use the new fixed exponent $L$ in
$\kappa=L+\Delta(a-1)$. The exponent-replacement congruence still uses
base $B$. There are twenty equations; all abbreviations are charged in
the certificate.

\begin{theorem}\label{thm:short118}
For every admissible index just defined, the shortened system represents
the same recursively enumerable set as the quadratic system from which
the index was constructed. It has a certificate of 118 instructions:
65 multiplications and 53 additions, followed by twenty equality tests.
\end{theorem}

\begin{proof}
We first verify the hypotheses needed to apply the Pell lemmas without
assuming the exponential equation. Equation~\eqref{eq:short-bound}
gives
\[
 e<q,\qquad \ell<q,\qquad C^4<q,
 \qquad (b-1)\ell<\ell C^4<q.
\]
Since $C>B>b\ge2$, it follows that $q>B$ and $g<q$. Equation
\eqref{eq:short-geometric} gives $\lambda>q$, so, writing $z=Z/2$,
\[
 0<D_0=z\lambda-e<z\lambda,
 \qquad 2C^4D_0<2z\lambda q<B\lambda q.
\]
The blocks therefore satisfy
\begin{gather*}
 0<g<q,\qquad 0\le q-1-(b-1)\ell<q,\\
 0<e+\ell q<q^2,\qquad
 0<(B-Z)\lambda<(B-1)\lambda=q^2-1,\\
 0<S_3<B\lambda(1+q^2)
      =\frac{B(q^4-1)}{B-1}<2q^4<q^5,\\
 0<(B-2)q<q^2<q^5.
\end{gather*}
Thus their respective widths are $N_1=q$, $N_2=q^2$, $N_3=q^5$;
their product is $n=q^8$. The equation for $r$ gives $r\ge n^2-1\ge n$,
and \eqref{eq:short-H} yields
$3<3L\le B<q\le n\le r$. The elementary Pell gap restores the
stronger inequality required by Lemma~\ref{lem:226}. The Pell lemmas
now give $b\pow2$, $q=B^L$, and the central-binomial divisibility.
Since $H$ is a power of two, the block radices are powers of two, and the
three separate no-carry conditions follow.

Apply \cite[Lemma~2.9]{J82} to $Y=e+\ell q$ with digit base $Z$ and
$m=k=2L$. The polynomial
\[
 F(B)=e_0(B)+\ell_0(B)B^L
\]
has degree at most $L+M<2L$, digits less than $Z$, and $F(Z)=V$.
The congruence in \eqref{eq:short-packed}, the bound $Y<q^2=B^{2L}$,
the mask $(B-Z)\lambda$, and \eqref{eq:short-H} supply the lemma's
hypotheses. Hence $Y=F(B)$. The supplied and intended $e,\ell$ all lie
in $[0,q)$, so quotient-and-remainder uniqueness gives $e=e_0(B)$ and
$\ell=\ell_0(B)$. The first mask identifies the witness digits of $g$
and bounds them by $b$.

To identify what the last mask tests, regard $B$ temporarily as a formal
base. The coefficient of $B^L$ in
\[
 C^4(e_0(B)-z\lambda)
\]
is exactly $R$ evaluated at the encoded tuple. Indeed, the monomial
exponent $w(I)$ in $C^4$ has coefficient $c_I$ times its monomial value;
the complementary digit of $e_0-z\lambda$ is $P_I$. This includes
$I=\boldsymbol0$, where that digit is $-z$, precisely the constant
coefficient in \eqref{eq:short-normalization}.

Every coefficient of $e_0-z\lambda$ has absolute value at most $z$.
Thus every coefficient of the product has absolute value less than
$z((\nu+2)b)^4<B/2$. Its degree is at most $2L-1+4M<3L$.
The offset $(B/2)\lambda(1+q^2)$ supplies the digit $B/2$ at every
position below $4L$, preventing carries throughout. The third no-carry
condition therefore tests, by \cite[Lemma~2.8]{J82}, whether the
coefficient $R$ is zero. The equivalence following
\eqref{eq:short-normalization} proves soundness.

Conversely, take a solution of the original quadratic system, set
$h=\sigma=1$, and choose a power-of-two $b$ greater than the input and
all sixty witnesses. Form the intended codes using $B=Hb^4$ and
$q=B^L$. Their bounds are
\[
 e_0(B)<ZB^{L-1},\qquad \ell_0(B)<2B^M,
 \qquad C<2bB^M.
\]
Because $B>2Z$, the first is less than $q/2$. Also
\[
 \ell C^4<32b^4B^{5M}
       <\tfrac12 B^{5M+1}\le\tfrac12q,
\]
using $H>64$ and $L=625M\ge5M+1$. Consequently the new
$\alpha=q-e-\ell C^4$ is positive. The packed polynomial $F$ has
nonnegative digits and positive degree, so $F(B)>F(Z)$ and
$t=(F(B)-F(Z))/(B-Z)>0$. The coordinate $\sigma=1$ ensures $g>0$.
All remaining witnesses follow from the same digit and Pell necessity
constructions with the new base and exponent.

For the count, replace the three instructions computing $e\ell$,
$(e\ell)C^2$ and addition of $\alpha$ by the three instructions
computing $\ell C^4$, addition of $e$, and addition of $\alpha$.
This exchanges one multiplication for one addition but preserves the
length. All packing formulas retain their previous instruction counts.
The power chain now consists only of $q^2,q^4,q^8$; the instructions
for $q^3$ and $q^{16}$ disappear. This reduces 120 to 118.

The exact verifier checks all twenty residuals. Its shared (E7)
calculation records the adapted identity
\[
 F_{7,\mathrm{new}}=F_7+(\lambda F_3-F_2)q(n^2-1),
\]
where $F_2,F_3$ now refer to
\eqref{eq:short-geometric}. Thus the symbolic check accounts for the
dependence on the preceding equations as well as the operation count.
\end{proof}

\begin{corollary}\label{thm:best114}
With the same admissible index $(Z,V,H)$, the shortened system together
with the signed Pell replacement of Proposition~\ref{lem:signed-pell} and the
linear ratio replacement of Proposition~\ref{lem:linear-ratio} has a
114-instruction certificate: 63 multiplications and 51 additions,
followed by 21 equality tests. It has 33 positive unknowns and the four
parameters $x,Z,V,H$, and defines the same recursively enumerable set.
\end{corollary}

\begin{proof}
Both Pell replacements apply to the shortened system: the relevant
equations retain their form, and the bounds $q>B>b\ge2$, $n=q^8\ge2$
and $a\ge20r$ supply their hypotheses. Their positive-domain equivalences
therefore compose with Theorem~\ref{thm:short118}. The final certificate
is independently checked by
\file{round5\_1980\_certificate.py}; its exact residual checks and
instruction histogram give the asserted 114 instructions and 21 tests.
The linear ratio replacement adds one positive unknown and one equation.
\end{proof}

For comparison, applying the three independent changes in this order gives:
\begin{center}
\begin{tabular}{@{}lrrr@{}}
\toprule
Construction & Multiplications & Additions & Total\\
\midrule
Packed encoding of Theorem~\ref{thm:120} & 68 & 52 & 120\\
Shorter masks & 65 & 53 & 118\\
Also the signed Pell parameter & 65 & 51 & 116\\
Also the positive interval & 63 & 51 & 114\\
\bottomrule
\end{tabular}
\end{center}

This is a further verified upper bound, not an assertion of minimality.
The choice $L=5^{64}$ also makes the fixed numeral inexpensive to
generate: after $2=1+1$, $4=2+2$, $5=4+1$, six successive squarings
give $5^2,5^4,\ldots,5^{64}$. Thus, if $x,Z,V,H$ remain supplied
parameters while $2,4,5,L$ must be generated from $1$, nine additional
instructions suffice, giving a total of 123.
